\documentclass[11pt]{article}

\usepackage{amsmath, amssymb, amsthm, amsfonts, mathtools}
\usepackage[hidelinks]{hyperref}
\usepackage{geometry}
\geometry{a4paper, margin=1in}

\theoremstyle{plain}
\newtheorem{theorem}{Theorem}[section]
\newtheorem{lemma}[theorem]{Lemma}
\newtheorem{corollary}[theorem]{Corollary}
\newtheorem{proposition}[theorem]{Proposition}

\theoremstyle{definition}
\newtheorem{definition}[theorem]{Definition}

\theoremstyle{remark}
\newtheorem{remark}[theorem]{Remark}
\newtheorem{example}[theorem]{Example}

\title{A Period Formula for Rational\\ Cut-and-Project Strip Projections\\ with Multiplicity}
\author{
    Dirk Kunert\thanks{Independent Researcher. Email: \texttt{dirk.kunert@gmail.com}}
}
\date{\today}

\begin{document}

\maketitle

\begin{abstract}
Consider the one-dimensional cut-and-project construction with coprime
positive integers \(\alpha,\beta\) defining the rational slope
\(\alpha/\beta\), and window parameter \(\omega \ge 0\).  Projecting lattice
points from the resulting strip onto the line and counting them with
multiplicity gives a multiset whose consecutive gaps form a periodic
sequence.  We prove that its minimal period is
\[
\lambda =
\begin{cases}
N, & D\nmid N,\\
N/D, & D\mid N,
\end{cases}
\]
where
\[
N=\lfloor\omega\alpha\rfloor+\lfloor\omega\beta\rfloor+1,
\qquad
D=\alpha^2+\beta^2.
\]
The proof reduces the geometry to the distribution of \(N\) consecutive
residue classes modulo~\(D\).  When this distribution is non-uniform, the
residue classes with larger multiplicity form a proper cyclic interval in
the natural residue parametrisation; its only translational symmetry is the
identity, which rules out any smaller period.  We also discuss how the
formula simplifies when multiplicities are discarded.  The algebraic core of
the proof has been machine-checked in Lean~4.
\end{abstract}

\medskip
\noindent\textbf{Mathematics Subject Classification (2020):}
Primary 52C23; Secondary 11B25, 37B10, 68V20, 11A07.

\medskip
\noindent\textbf{Keywords:}
cut-and-project sets, projected multisets, gap sequences, period length,
rational slope, residue classes, lattice projection, Lean~4 formalisation.

% =====================================================================
\section{Introduction}\label{sec:intro}
% =====================================================================

Cut-and-project constructions provide a standard way of producing ordered
point sets from lattices.  In the one-dimensional planar case, one selects
lattice points lying in a strip around a line and projects them onto that
line.  When the slope is irrational, the resulting objects are typically
aperiodic and belong to the classical theory of model sets and aperiodic
order.  When the slope is rational, the resulting structure is periodic and
is therefore often regarded as a degenerate or elementary case from the
aperiodic-order point of view.

This viewpoint goes back in part to Yves Meyer's foundational work on
harmonious sets and model sets~\cite{Meyer1972}.  Meyer's work became one
of the mathematical foundations of the modern theory of quasicrystals and
aperiodic order, and was recognised by the 2017 Abel Prize.

Nevertheless, even in the rational case there are natural quantitative
questions.  One such question is the following: after projecting the accepted
lattice points and ordering the projected values along the line, what is the
minimal period of the sequence of consecutive gaps?  This question depends
on a convention.  If projected values are counted with multiplicity, then
distinct lattice points projecting to the same value give rise to zero gaps.
If multiplicities are discarded, the resulting ordinary point-set gap
sequence may have a smaller period.  The main theorem of this paper is
stated for the projected multiset.

Let the rational slope be \(\alpha/\beta\), with
\(\gcd(\alpha,\beta)=1\), and let \(\omega\ge0\) be the window parameter.
We prove that the multiset gap sequence has minimal period
\[
    \lambda =
    \begin{cases}
        N, & D\nmid N,\\[3pt]
        N/D, & D\mid N,
    \end{cases}
\]
where
\[
    N=\lfloor\omega\alpha\rfloor+\lfloor\omega\beta\rfloor+1,
    \qquad
    D=\alpha^2+\beta^2.
\]
Here \(N\) is the number of admissible values of the invariant
\(r=\alpha x-\beta y\), and \(D\) is the common difference of the arithmetic
progressions obtained after projection.

The proof is elementary.  The accepted lattice points decompose into
parallel tracks indexed by the admissible values of \(r\).  Each track gives
an arithmetic progression of projected values with common difference \(D\),
and the relevant residues modulo \(D\) are obtained from \(N\) consecutive
integers by multiplication with a unit modulo \(D\).  If \(D\nmid N\), the
residue classes with larger multiplicity form a proper cyclic interval in
the natural residue parametrisation; its only translational symmetry is the
identity, which rules out any smaller period.  If \(D\mid N\), the residues
are uniformly distributed and the gap sequence collapses to a simple pattern
of period \(N/D\).

The algebraic core of the proof has been formalised in Lean~4.  The
formalisation verifies the residue-class and minimal-period arguments; the
geometric construction is represented at the combinatorial level described
above.

The paper is organised as follows.  Section~\ref{sec:setup} gives the
construction and fixes the multiplicity convention.
Section~\ref{sec:mainresult} states the main result.
Section~\ref{sec:proof} proves the generic case \(D \nmid N\).
Section~\ref{sec:degenerate} treats the uniform case \(D \mid N\).
Section~\ref{sec:corollaries} records some immediate corollaries.
Section~\ref{sec:lean} describes the Lean~4 formalisation.
Section~\ref{sec:literature} discusses related literature.
Section~\ref{sec:discussion} discusses the set-valued variant and open
questions.

% =====================================================================
\section{Construction of the Sequences}\label{sec:setup}
% =====================================================================

\subsection{Accepted lattice points}\label{subsec:construction}

Let \(\alpha, \beta \in \mathbb{N}\) with \(\gcd(\alpha, \beta) = 1\) and let
\(\omega \in \mathbb{R}_{\ge 0}\).
Consider the line \(f(x) = \frac{\alpha}{\beta}\,x\) in \(\mathbb{R}^2\)
together with the strip bounded by
\[
    l(x) = \frac{\alpha}{\beta}(x - \omega)
    \qquad\text{and}\qquad
    u(x) = \frac{\alpha}{\beta}\,x + \omega.
\]
The set of \emph{accepted lattice points} is
\begin{equation}\label{eq:points}
    \mathbf{C} \;=\; \left\{
    \begin{pmatrix} x \\ y \end{pmatrix} \in \mathbb{Z}^2
    \;\middle|\;
    y \in
    \left[\frac{\alpha}{\beta}(x - \omega),\;
          \frac{\alpha}{\beta}\,x + \omega\right]
    \right\}.
\end{equation}

\subsection{Projection and ordering}\label{subsec:projection}

The orthogonal projection of \((x, y) \in \mathbb{R}^2\) onto the line \(f\)
is
\begin{equation}\label{eq:projection}
    \mathbf{p} \;=\;
    \frac{1}{\alpha^2 + \beta^2}
    \begin{pmatrix}
        \beta^2 & \alpha\beta \\
        \alpha\beta & \alpha^2
    \end{pmatrix}
    \begin{pmatrix} x \\ y \end{pmatrix}.
\end{equation}
The \(x\)-coordinate of the projected point is
\[
    p_x = \frac{\beta(\beta x + \alpha y)}{\alpha^2 + \beta^2}.
\]
Since the factor \(\frac{\beta}{\alpha^2+\beta^2}\) is a positive constant,
the ordering of projected points is determined by
\begin{equation}\label{eq:ptilde}
    \tilde{p} \;:=\; \beta x + \alpha y.
\end{equation}

\subsection{Multiplicity convention}\label{subsec:multiplicity}

Throughout this paper, projected values are counted with multiplicity.
Thus, if two distinct accepted lattice points have the same projected
coordinate, they contribute two consecutive entries to the sorted projected
sequence and may therefore produce a zero gap.

Equivalently, we study the gap sequence of the projected lattice-point
\emph{multiset}, not only the underlying projected point set.  If
multiplicities are discarded, the resulting ordinary point-set gap sequence
may have a smaller period.  In particular, in the case \(D \mid N\), the
set-valued projected sequence has constant gap and period~\(1\), whereas
the multiset sequence has period~\(N/D\).

\subsection{Distance sequence and period}\label{subsec:distance}

Sort the multiset of all \(\tilde{p}\)-values in nondecreasing order, with
multiplicities retained, to obtain
\((\tilde{p}^{(i)}_s)_{i \in \mathbb{Z}}\), and define the
\emph{difference sequence}
\begin{equation}\label{eq:delta}
    \delta^{(i)} \;:=\; \tilde{p}^{(i+1)}_s - \tilde{p}^{(i)}_s.
\end{equation}
Since orthogonal projection onto the line rescales the ordered
\(\tilde p\)-coordinate by a positive constant, the period length of
\((\delta^{(i)})\) equals the period length of the corresponding Euclidean
distance sequence between consecutive projected multiset entries, including
zero distances arising from coincident projected points.

\begin{definition}\label{def:period}
The sequence \((\delta^{(i)})_{i \in \mathbb{Z}}\) has \emph{period length}
\(\lambda \in \mathbb{N}\) if \(\lambda\) is the smallest positive integer
such that \(\delta^{(i+\lambda)} = \delta^{(i)}\) for all~\(i\).
\end{definition}

\subsection{The trivial case \(\omega = 0\)}\label{subsec:omega0}

When \(\omega = 0\), the only accepted lattice points lie on \(f\) itself:
\((x, y) = k(\beta, \alpha)\) for \(k \in \mathbb{Z}\).
Then
\[
    \tilde{p} = \beta \cdot k\beta + \alpha \cdot k\alpha
    = k(\alpha^2+\beta^2),
\]
so all differences equal \(\alpha^2+\beta^2\) and \(\lambda = 1\).
Note that the formula~\eqref{eq:formula} agrees: \(N = 0 + 0 + 1 = 1\)
and \(D = \alpha^2+\beta^2 \ge 2\), so \(D \nmid 1\) and \(\lambda = N = 1\).

% =====================================================================
\section{Main Result}\label{sec:mainresult}
% =====================================================================

\begin{theorem}[Period length formula]\label{thm:main}
Let \(\alpha, \beta \in \mathbb{N}\) with \(\gcd(\alpha, \beta) = 1\) and let
\(\omega \in \mathbb{R}_{\ge 0}\).  Define
\begin{equation}\label{eq:ND}
    N \;:=\; \lfloor\omega\alpha\rfloor + \lfloor\omega\beta\rfloor + 1,
    \qquad
    D \;:=\; \alpha^2 + \beta^2.
\end{equation}
Then the period length of the difference sequence~\eqref{eq:delta} is
\begin{equation}\label{eq:formula}
    \lambda_{\alpha,\beta} \;=\;
    \begin{cases}
        N & \text{if } D \nmid N, \\[3pt]
        N / D & \text{if } D \mid N.
    \end{cases}
\end{equation}
\end{theorem}

\begin{remark}[Interpretation of \(N\)]\label{rem:N}
The quantity \(N\) equals the number of integers in the interval
\([-\beta\omega,\; \alpha\omega]\), which is the range of the invariant
\(r = \alpha x - \beta y\) (see Section~\ref{subsec:step1}).
Geometrically, \(N\) counts the number of tracks, namely the parallel
arithmetic progressions of projected values counted with multiplicity.
\end{remark}

\begin{remark}[Interpretation of \(D\)]\label{rem:D}
The quantity \(D = \alpha^2+\beta^2\) is the common difference of each
arithmetic progression of \(\tilde{p}\)-values.
It also equals \(\lVert(\beta,\alpha)\rVert^2\), the squared length of the
direction vector of~\(f\).
\end{remark}

\begin{remark}[Computational verification]\label{rem:verification}
The formula has been verified against:
\begin{itemize}
    \item 51{,}012 test cases from~\cite{Kunert2025} with
          \(\alpha, \beta\) up to \({\sim}100{,}000\) and rational \(\omega\)
          up to \({\sim}10\), all satisfying \(D \nmid N\)
          (51{,}011 with \(N < D\), one with \(N \ge D\));
    \item 468 deliberately constructed cases in the degenerate regime
          \(D \mid N\) (with \(\alpha, \beta \le 10\) and \(N\) up to
          \({\sim}500\)), verifying the \(\lambda = N/D\) branch.
\end{itemize}
All 51{,}480 cases produce exact matches with zero failures.
\end{remark}

% =====================================================================
\section{Proof: The Generic Case \(D \nmid N\)}\label{sec:proof}
% =====================================================================

\subsection{Step 1: Range of the invariant}\label{subsec:step1}

The strip condition~\eqref{eq:points}, multiplied by \(\beta > 0\), reads
\[
    \alpha(x - \omega) \;\le\; \beta y \;\le\; \alpha x + \beta\omega,
\]
equivalently
\begin{equation}\label{eq:r_range}
    -\beta\omega \;\le\;
    \underbrace{\alpha x - \beta y}_{=:\,r}
    \;\le\; \alpha\omega.
\end{equation}
Since \(r \in \mathbb{Z}\), it takes values in
\begin{equation}\label{eq:R}
    \mathcal{R} \;:=\;
    \bigl\{-\lfloor\beta\omega\rfloor,\;
    -\lfloor\beta\omega\rfloor + 1,\; \ldots,\;
    \lfloor\alpha\omega\rfloor\bigr\},
\end{equation}
using the identity \(\lceil{-x}\rceil = -\lfloor x \rfloor\).
Hence
\[
    |\mathcal{R}| =
    \lfloor\alpha\omega\rfloor + \lfloor\beta\omega\rfloor + 1 = N.
\]

The invariant \(r\) is constant along each track: if \((x, y)\) and
\((x + \beta, y + \alpha)\) are both accepted, they share the same~\(r\),
since
\[
    \alpha(x+\beta) - \beta(y+\alpha) = \alpha x - \beta y.
\]
The step \((\beta, \alpha)\) is the direction of~\(f\).

\subsection{Step 2: Each value of \(r\) yields one arithmetic progression}
\label{subsec:step2}

For a fixed \(r \in \mathcal{R}\), the system
\begin{equation}\label{eq:system}
    \beta x + \alpha y = \tilde{p}, \qquad
    \alpha x - \beta y = r
\end{equation}
has determinant \(-(\alpha^2+\beta^2) = -D \ne 0\), giving the unique solution
\begin{equation}\label{eq:solution}
    x = \frac{\beta\tilde{p} + \alpha r}{D}, \qquad
    y = \frac{\alpha\tilde{p} - \beta r}{D}.
\end{equation}
Integrality requires \(D \mid (\beta\tilde{p} + \alpha r)\), i.e.
\begin{equation}\label{eq:residue}
    \tilde{p} \;\equiv\; c_r \;:=\; -\alpha\beta^{-1}r \pmod{D},
\end{equation}
where \(\beta^{-1}\) is the modular inverse of \(\beta\) modulo \(D\)
(existence: Corollary~\ref{cor:inverse}).
The set of \(\tilde{p}\)-values for invariant~\(r\) is the arithmetic
progression
\begin{equation}\label{eq:Pr}
    P_r \;=\; \{c_r + kD : k \in \mathbb{Z}\}.
\end{equation}
The complete multiset of projected values is
\[
    S = \bigcup_{r \in \mathcal{R}} P_r,
\]
where the union is understood as a multiset union.

\subsection{Step 3: Coprimality lemmas}\label{subsec:step3}

\begin{lemma}\label{lem:gcd}
\(\gcd(\alpha, D) = \gcd(\beta, D) = 1\).
\end{lemma}

\begin{proof}
\(D = \alpha^2 + \beta^2 \equiv \beta^2 \pmod{\alpha}\), so any common
divisor of \(\alpha\) and \(D\) divides \(\beta^2\).  But
\(\gcd(\alpha,\beta)=1\) implies \(\gcd(\alpha,\beta^2)=1\).
The argument for \(\beta\) is symmetric.
\end{proof}

\begin{corollary}\label{cor:inverse}
The modular inverse \(\beta^{-1} \pmod{D}\) exists, and the map
\(r \mapsto c_r = -\alpha\beta^{-1}r \bmod D\) is a bijection on
\(\mathbb{Z}/D\mathbb{Z}\).
\end{corollary}

\subsection{Step 4: Structure of residues}\label{subsec:step4}

The \(N\) residues \(\{c_r : r \in \mathcal{R}\}\) are the image of \(N\)
consecutive integers under multiplication by
\(m := -\alpha\beta^{-1} \bmod D\), with \(\gcd(m, D) = 1\).  Explicitly:
\begin{equation}\label{eq:residues}
    \{c_r\}_{r \in \mathcal{R}} \;=\;
    \{m \cdot r_0,\; m(r_0+1),\; \ldots,\; m(r_0+N-1)\} \bmod D,
\end{equation}
where \(r_0 = -\lfloor\beta\omega\rfloor\).

\begin{lemma}[Non-uniform residue distribution]\label{lem:nonuniform}
If \(D \nmid N\), the multiset of \(N\) residues
\(\{c_r \bmod D : r \in \mathcal{R}\}\) is not uniformly distributed over
\(\mathbb{Z}/D\mathbb{Z}\).  More precisely, writing \(N = qD + s\) with
\(0 < s < D\), exactly \(s\) residue classes are hit \(q+1\) times and
\(D - s\) residue classes are hit \(q\) times.
\end{lemma}

\begin{proof}
The map \(r \mapsto mr \bmod D\) is a bijection on
\(\mathbb{Z}/D\mathbb{Z}\) by Corollary~\ref{cor:inverse}.
As \(r\) runs through \(N\) consecutive integers, the residues \(mr \bmod D\)
cycle through all \(D\) classes with period~\(D\).
After \(q\) complete cycles, all classes have been hit exactly \(q\) times.
The remaining \(s = N - qD\) values hit \(s\) additional classes exactly once
each.
Hence \(s\) classes have multiplicity \(q + 1\) and \(D - s\) classes have
multiplicity~\(q\).
\end{proof}

\subsection{Step 5: Period equals \(N\) when \(D \nmid N\)}
\label{subsec:step5}

\begin{proof}
Within each interval \([kD,\, (k+1)D)\) on the \(\tilde{p}\)-axis, exactly
\(N\) points from the multiset \(S\) occur, counted with multiplicity.
Their positions within the interval are determined by the residues
\(c_r \bmod D\).

The sorted difference sequence therefore has a repeating block of length
\(N\): after \(N\) entries, the sorted multiset is translated by \(D\), so
\[
    \tilde p_s^{(i+N)} = \tilde p_s^{(i)} + D
\]
for all \(i\).  Hence
\[
    \delta^{(i+N)}=\delta^{(i)}
\]
for all \(i\), and therefore \(\lambda \mid N\).

\medskip
\noindent\textbf{Minimality.}
We now prove that no proper divisor of \(N\) is a period, using a
multiplicity-function argument on \(\mathbb{Z}/D\mathbb{Z}\).

Define the multiplicity function
\[
    \mu \colon \mathbb{Z}/D\mathbb{Z} \to \mathbb{N}_0
\]
by
\[
    \mu(c) \;:=\; \bigl|\{r \in \mathcal{R} : c_r \equiv c \pmod{D}\}\bigr|.
\]
By Lemma~\ref{lem:nonuniform}, with \(N = qD + s\) and \(0 < s < D\),
the function \(\mu\) takes exactly two values: \(q+1\) on a set of \(s\)
residue classes and \(q\) on the remaining \(D - s\) classes.

It is convenient to pass to the \(m\)-stepping parametrisation
\[
    f(x) := \mu(mx), \qquad x \in \mathbb{Z}/D\mathbb{Z},
\]
where \(m=-\alpha\beta^{-1}\pmod D\).
Since \(m\) is invertible modulo \(D\), this is only a relabelling of the
same data.  In this parametrisation, the heavy set
\[
    H := \{x \in \mathbb{Z}/D\mathbb{Z} : f(x)=q+1\}
\]
is a contiguous cyclic interval of length \(s\):
\[
    H = \{x_0, x_0+1, \ldots, x_0+s-1\}\pmod D
\]
for some \(x_0\), because the final \(s\) elements of the consecutive block
\(r_0,r_0+1,\ldots,r_0+N-1\) contribute one extra hit exactly on \(s\)
consecutive values in the \(m\)-stepping order.

Now suppose, for contradiction, that \(\lambda' < N\) is a period of the gap
sequence, and write \(N = k\lambda'\) with \(k>1\).
Let
\[
    \sigma := \delta^{(0)} + \delta^{(1)} + \cdots + \delta^{(\lambda'-1)}.
\]
Since the gap sequence has period \(\lambda'\), we have
\[
    \tilde p_s^{(i+\lambda')} = \tilde p_s^{(i)} + \sigma
\]
for all \(i\), so the entire multiset \(S\) is invariant under translation by
\(\sigma\).  On the other hand, \(N\) consecutive gaps always sum to \(D\).
Hence
\[
    k\sigma = D,
\]
so \(\sigma=D/k\) is a positive integer and \(k \mid D\).

The translation invariance of \(S\) implies invariance of the residue
multiplicity function:
\[
    \mu(c+\sigma)=\mu(c) \qquad \text{for all } c\in \mathbb{Z}/D\mathbb{Z}.
\]
Equivalently, the relabelled function \(f\) satisfies
\[
    f(x+\tau)=f(x), \qquad \tau := m^{-1}\sigma \not\equiv 0 \pmod D.
\]
Therefore its heavy set \(H\) must satisfy \(H+\tau = H\).

But a proper nonempty cyclic interval in \(\mathbb{Z}/D\mathbb{Z}\) has
trivial stabiliser: if \(0<s<D\) and
\[
    H=\{x_0,x_0+1,\ldots,x_0+s-1\}\pmod D,
\]
then \(H+\tau=H\) forces \(\tau\equiv 0\pmod D\).
Indeed, for \(1\le \tau \le D-1\), the translate \(H+\tau\) has a different
boundary and cannot coincide with \(H\) unless \(H\) is either empty or all of
\(\mathbb{Z}/D\mathbb{Z}\), neither of which occurs because \(0<s<D\).

This contradiction shows that no proper divisor \(\lambda' < N\) can be a
period.  Since we already know \(\lambda \mid N\), it follows that
\(\lambda=N\).
\end{proof}

% =====================================================================
\section{The Degenerate Case \(D \mid N\)}\label{sec:degenerate}
% =====================================================================

\begin{theorem}\label{thm:degenerate}
If \(D \mid N\), then \(\lambda = N/D\).
\end{theorem}

\begin{proof}
When \(D \mid N\), the \(N\) values of \(r\) span exactly \(N/D\) complete
cycles of the map \(r \mapsto mr \bmod D\).
Since this map is a bijection on \(\mathbb{Z}/D\mathbb{Z}\), each residue
class \(c \in \{0, 1, \ldots, D-1\}\) is hit exactly \(N/D\) times.
Consequently, each residue class contributes \(N/D\) projected multiset
entries per interval of length~\(D\) on the \(\tilde{p}\)-axis.  By the
uniform distribution, every integer \(\tilde{p}\) is achieved with
multiplicity~\(N/D\).

The sorted sequence of \(\tilde{p}\)-values, counted with multiplicity, thus
consists of each integer repeated \(N/D\) times:
\[
    \ldots,\;
    \underbrace{t,\, t,\, \ldots,\, t}_{N/D},\;
    \underbrace{t+1,\, t+1,\, \ldots,\, t+1}_{N/D},\;
    \ldots
\]
The difference sequence is
\[
    \underbrace{0,\, 0,\, \ldots,\, 0}_{N/D - 1},\; 1,\;
    \underbrace{0,\, 0,\, \ldots,\, 0}_{N/D - 1},\; 1,\;
    \ldots
\]
which has period \(N/D\).
\end{proof}

\begin{example}\label{ex:degenerate}
For \(\alpha = 2\), \(\beta = 1\), \(\omega = 3\):
\[
    N = \lfloor 6 \rfloor + \lfloor 3 \rfloor + 1 = 10,
    \qquad
    D = 5.
\]
Since \(5 \mid 10\), the formula gives \(\lambda = 10/5 = 2\).
Direct computation confirms the difference sequence
\((0, 1, 0, 1, \ldots)\) with period~\(2\).
\end{example}

\begin{example}\label{ex:degenerate2}
For \(\alpha = 3\), \(\beta = 2\), \(\omega = 5\):
\[
    N = \lfloor 15 \rfloor + \lfloor 10 \rfloor + 1 = 26,
    \qquad
    D = 13.
\]
Since \(13 \mid 26\), the formula gives \(\lambda = 26/13 = 2\).
\end{example}

\begin{remark}
The boundary case \(N = D\), i.e. \(N/D = 1\), gives \(\lambda = 1\):
every integer is a projected value with multiplicity~\(1\), and all gaps
equal~\(1\).
\end{remark}

% =====================================================================
\section{Corollaries}\label{sec:corollaries}
% =====================================================================

Define \(\Lambda_{\alpha,\beta} := \alpha + \beta + 1\).

\begin{corollary}[Finiteness]\label{cor:finite}
For all \(\alpha, \beta \in \mathbb{N}\) with \(\gcd(\alpha,\beta)=1\) and
all \(\omega \in \mathbb{R}_{\ge 0}\), the period length~\(\lambda\) is
finite.
\end{corollary}

\begin{proof}
\(N\) is a finite positive integer, and \(\lambda \le N\).
\end{proof}

\begin{corollary}[Period at \(\omega = 1\)]\label{cor:omega1}
For \(\omega = 1\), \(\lambda = \alpha + \beta + 1 = \Lambda_{\alpha,\beta}\).
\end{corollary}

\begin{proof}
\(N = \alpha + \beta + 1\) and \(D = \alpha^2+\beta^2\).
For \((\alpha,\beta) \ne (1,1)\), we have
\[
    \alpha+\beta+1 < \alpha^2+\beta^2,
\]
so \(D \nmid N\) and \(\lambda = N = \alpha+\beta+1\).
For \((\alpha,\beta) = (1,1)\), \(N = 3\), \(D = 2\), and again
\(D \nmid N\), so \(\lambda = 3 = \Lambda_{1,1}\).
\end{proof}

\begin{corollary}[\(0 < \omega < 1\)]\label{cor:small_omega}
If \(0 < \omega < 1\), then
\[
    \lambda < \Lambda_{\alpha,\beta}.
\]
\end{corollary}

\begin{proof}
\(\lfloor\omega\alpha\rfloor \le \alpha - 1\) and
\(\lfloor\omega\beta\rfloor \le \beta - 1\), so
\[
    N \le \alpha + \beta - 1 < \Lambda_{\alpha,\beta}.
\]
Since \(\alpha + \beta - 1 < \alpha^2 + \beta^2 = D\) for all
\(\alpha, \beta \ge 1\), we have \(N < D\), hence \(D \nmid N\) and
\[
    \lambda = N < \Lambda_{\alpha,\beta}.
\]
\end{proof}

\begin{corollary}[\(1 < \omega < 2\)]\label{cor:medium_omega}
If \(1 < \omega < 2\) and \(D \nmid N\), then
\[
    \lambda \ge \Lambda_{\alpha,\beta}.
\]
\end{corollary}

\begin{proof}
\(\lfloor\omega\alpha\rfloor \ge \alpha\) and
\(\lfloor\omega\beta\rfloor \ge \beta\), so
\[
    N \ge \alpha + \beta + 1 = \Lambda_{\alpha,\beta}.
\]
Since \(D \nmid N\), we have \(\lambda = N \ge \Lambda_{\alpha,\beta}\).
\end{proof}

\begin{remark}
The condition \(D \nmid N\) in Corollary~\ref{cor:medium_omega} cannot be
dropped.  For \(\alpha = 2\), \(\beta = 1\), \(\omega = 3/2\), one has
\[
    N = \lfloor 3 \rfloor + \lfloor 3/2 \rfloor + 1 = 5 = D.
\]
Since \(D \mid N\), we get \(\lambda = N/D = 1\), which is less than
\(\Lambda_{2,1} = 4\).
\end{remark}

\begin{corollary}[\(\omega > 2\)]\label{cor:large_omega}
If \(\omega > 2\) and \(D \nmid N\), then
\[
    \lambda > \Lambda_{\alpha,\beta}.
\]
\end{corollary}

\begin{proof}
\(\lfloor\omega\alpha\rfloor \ge 2\alpha\) and
\(\lfloor\omega\beta\rfloor \ge 2\beta\), so
\[
    N \ge 2\alpha + 2\beta + 1 > \Lambda_{\alpha,\beta}.
\]
Since \(D \nmid N\), \(\lambda = N > \Lambda_{\alpha,\beta}\).
\end{proof}

\begin{remark}
In both corollaries above, the degenerate case \(D \mid N\) can produce
small periods: \(\lambda = N/D\) may be much less than
\(\Lambda_{\alpha,\beta}\).  This occurs when the window width is tuned so
that the number of tracks is an exact multiple of \(D = \alpha^2+\beta^2\).
\end{remark}

% =====================================================================
\section{Machine-Checked Formalisation}\label{sec:lean}
% =====================================================================

The algebraic core of the proof has been formalised in
Lean~4~\cite{Moura2021} using the Mathlib library~\cite{Mathlib2020}.
The formalisation comprises approximately 1{,}300 lines of verified code
with no \texttt{sorry} axioms, and is available in the companion
repository~\cite{Kunert2025} under \texttt{Lean/CutAndProject/}.

\subsection{Architecture}\label{subsec:lean_arch}

The proof is structured in two layers.  The first layer proves
Theorem~\ref{thm:main} from four abstract axioms collected in a
typeclass \texttt{GeometricProjection}:
\begin{enumerate}
    \item \(N > 0\) \hfill (\texttt{N\_pos})
    \item The difference sequence has period \(N\)
          \hfill (\texttt{period\_N})
    \item If \(D \mid N\), the minimal period is \(N/D\)
          \hfill (\texttt{period\_degenerate})
    \item Every period \(L\) induces a residue-preserving translation
          by some \(\sigma\) with \(\sigma N = LD\)
          \hfill (\texttt{sigma\_of\_period})
\end{enumerate}
Given these axioms, the generic minimality argument
(\texttt{generic\_minimality}, corresponding to
Section~\ref{subsec:step5}) and the case dispatch
(\texttt{main\_theorem}, corresponding to Theorem~\ref{thm:main})
are proved purely algebraically.

The second layer constructs a concrete instance of this typeclass
by defining the sorted multiset and difference sequence explicitly.
Each axiom is discharged by a separate lemma
(\texttt{N\_pos\_concrete}, \texttt{period\_N\_concrete},
\texttt{period\_degenerate\_concrete},
\texttt{sigma\_of\_period\_concrete}).

\subsection{Correspondence with the paper}\label{subsec:lean_corr}

Table~\ref{tab:lean} gives the correspondence between the main
results of this paper and the Lean~4 declarations in
\texttt{Basic.lean}.

\begin{table}[ht]
\centering
\small
\begin{tabular}{lll}
\hline
\textbf{Paper} & \textbf{Lean~4 name} & \textbf{Line} \\
\hline
Lemma~\ref{lem:gcd} (\(\gcd(\alpha,D)=1\))
  & \texttt{coprime\_alpha\_D} & 11 \\
Lemma~\ref{lem:gcd} (\(\gcd(\beta,D)=1\))
  & \texttt{coprime\_beta\_D} & 20 \\
Corollary~\ref{cor:inverse} (residue bijection)
  & \texttt{residue\_bijection} & 43 \\
Lemma~\ref{lem:nonuniform} (non-uniform distribution)
  & \texttt{non\_uniform\_residue\_distribution} & 140 \\
Uniform distribution (\(D \mid N\))
  & \texttt{uniform\_residue\_distribution} & 215 \\
Heavy set is cyclic interval
  & \texttt{heavy\_set\_is\_cyclic\_interval} & 269 \\
Trivial stabiliser
  & \texttt{cyclic\_interval\_stabilizer\_trivial} & 334 \\
Section~\ref{subsec:step5} (generic minimality)
  & \texttt{generic\_minimality} & 388 \\
Theorem~\ref{thm:degenerate} (\(D \mid N\) case)
  & \texttt{period\_degenerate\_concrete} & 1233 \\
\(\sigma N = LD\) from periodicity
  & \texttt{sigma\_of\_period\_concrete} & 1206 \\
Theorem~\ref{thm:main} (main result)
  & \texttt{main\_theorem} & 1283 \\
\hline
\end{tabular}
\caption{Correspondence between paper results and Lean~4 declarations.}
\label{tab:lean}
\end{table}

\subsection{Scope of the formalisation}\label{subsec:lean_scope}

The formalised proof covers the algebraic and combinatorial core:
coprimality, residue distribution, cyclic interval structure, the trivial
stabiliser lemma, the generic minimality argument, the degenerate case, and
the final case dispatch.

The geometric construction, namely defining the cut-and-project set,
projecting lattice points, and sorting the multiset of \(\tilde{p}\)-values,
is modelled abstractly via the \texttt{GeometricProjection} typeclass rather
than formalised from first principles.  Concretely, the sorted multiset is
defined combinatorially as
\[
    \tilde{p}^{(i)}_s := V(i \bmod N) + \lfloor i/N \rfloor \cdot D,
\]
where \(V\) is the quantile function of the cumulative residue distribution.
The bridge lemma \texttt{count\_hits\_eq\_sorted\_count} verifies that this
construction is consistent with the residue counting function.

% =====================================================================
\section{Related Literature}\label{sec:literature}
% =====================================================================

\subsection{The three-distance theorem}

The three-distance theorem of Steinhaus, S\'{o}s, Sur\'{a}nyi, and
\'{S}wierczkowski states that placing \(n\) points \(\{k\theta\}\) on a unit
circle produces at most three distinct gap lengths~\cite{ThreeGap,
BertheReutenauer2024}.  More recent treatments include the lattice proof by
Marklof and Str\"{o}mbergsson~\cite{MarklofStrombergsson2017}.

This theorem is related in spirit but differs from Theorem~\ref{thm:main} in
two essential ways:
\begin{itemize}
    \item It concerns rotations on a circle rather than lattice projections
          through a strip of finite width.
    \item It characterises the number of distinct gap lengths, not the
          minimal period length of a gap sequence.
\end{itemize}

\subsection{Christoffel words and Sturmian sequences}

The lower Christoffel word of slope \(p/q\) with \(\gcd(p,q)=1\) is a binary
word of length \(p+q\), and the mechanical Sturmian sequence of rational
slope \(p/q\) has period \(p+q\)~\cite{BerstelEtAl2008}.

The difference from Theorem~\ref{thm:main} is significant:
\begin{enumerate}
    \item Christoffel words have period \(p+q\), whereas at \(\omega = 1\)
          our multiset construction gives \(\alpha+\beta+1\).
    \item The Sturmian framework describes symbolic codings, not the
          projected multiset gap sequence considered here.
    \item The present result concerns distances between consecutive projected
          multiset entries, including zero gaps when projected values
          coincide.
\end{enumerate}

\subsection{Cut-and-project sets and model sets}

The mainstream literature on model sets~\cite{Senechal2009, Meyer1972,
HaynesEtAl2016} focuses largely on irrational slopes and aperiodic order.
Rational slopes yield periodic structures and are therefore often less
central from that perspective.

The author is not aware of a reference that states the above formula for the
minimal period of the multiset gap sequence associated with a rational
one-dimensional cut-and-project strip.  The result is elementary once the
appropriate residue-class description is introduced, but it appears useful to
record explicitly.

% =====================================================================
\section{Discussion}\label{sec:discussion}
% =====================================================================

\subsection{Why the formula splits into two cases}

The key structural distinction is whether the \(N\) tracks, viewed as
arithmetic progressions modulo~\(D\), distribute uniformly or not:
\begin{itemize}
    \item When \(D \nmid N\), the tracks distribute non-uniformly among the
          \(D\) residue classes.  The heavy residue classes form a proper
          cyclic interval in the natural residue parametrisation, whose only
          translational symmetry is the identity.  This gives
          \(\lambda = N\).
    \item When \(D \mid N\), each residue class is hit exactly \(N/D\) times.
          Every integer is a projected value with multiplicity \(N/D\),
          collapsing the gap structure to a simple pattern of period \(N/D\).
\end{itemize}

\subsection{Discarding multiplicities}\label{subsec:discarding}

The main theorem is stated for the projected multiset.  If projected points
are instead identified, the period may decrease.  The clearest example is
the degenerate case \(D \mid N\).  There every integer \(\tilde p\) occurs
with multiplicity \(N/D\).  After discarding multiplicities, the projected
set is simply \(\mathbb{Z}\) in the \(\tilde p\)-coordinate, so the ordinary
set-valued gap sequence has constant gap and period~\(1\).

Thus the branch \(\lambda=N/D\) is a genuine multiset phenomenon.  A complete
formula for the set-valued gap sequence, with multiplicities discarded in
all cases, is a natural companion question.

\subsection{Open problems}

\begin{enumerate}
    \item \textbf{Dependence on the window position.}
    The present paper considers a strip in the specific position used in
    \eqref{eq:points}.  If one translates the window transversally, the set of
    admissible invariant values changes by a shifted interval, but boundary
    effects may alter the interleaving pattern.  It would be interesting to
    determine whether the same period formula persists for arbitrary window
    offsets.

    \item \textbf{Set-valued gap sequence.}
    Determine the exact minimal period when projected points are counted
    without multiplicity.  The degenerate case \(D \mid N\) collapses to
    period~\(1\), but the behaviour in the non-uniform case should be
    described explicitly.

    \item \textbf{Higher dimensions.}
    The cut-and-project construction generalises to higher-dimensional
    lattices.  Whether analogous period formulas hold in dimension
    \(d \ge 2\) is unexplored.

    \item \textbf{Distribution of gap lengths.}
    For fixed \(\alpha, \beta, \omega\), one may ask for a precise
    characterisation of the distinct gap values and their multiplicities
    within one minimal period.
\end{enumerate}

% =====================================================================
\section*{Acknowledgements}
% =====================================================================

The author used ChatGPT, Gemini, and Claude as AI-assisted tools during
mathematical exploration, drafting, informal proof checking, and the
preparation of the Lean formalisation.  The author has independently checked
the mathematical content and accepts full responsibility for the results and
any errors.

% =====================================================================
\begin{thebibliography}{99}
% =====================================================================

\bibitem{BerstelEtAl2008}
Jean Berstel, Aaron Lauve, Christophe Reutenauer, and Franco V.\ Saliola,
\emph{Combinatorics on Words: Christoffel Words and Repetitions in Words},
CRM Monograph Series, Vol.~27,
American Mathematical Society, 2008.

\bibitem{BertheReutenauer2024}
Val\'{e}rie Berth\'{e} and Christophe Reutenauer,
``On the three distance theorem,''
\emph{The Mathematical Intelligencer}, 2024.
\textsc{doi}: \href{https://doi.org/10.1007/s00283-023-10316-z}%
{10.1007/s00283-023-10316-z}.

\bibitem{HaynesEtAl2016}
Alan Haynes, Henna Koivusalo, James Walton, and Lorenzo Sadun,
``Gaps problems and frequencies of patches in cut and project sets,''
\emph{Mathematical Proceedings of the Cambridge Philosophical Society},
vol.~161, no.~1, pp.~65--85, 2016.

\bibitem{Kunert2025}
Dirk Kunert,
\emph{Repository cut-and-project}, 2025.
Available at: \url{https://github.com/dkunert/cut-and-project}.

\bibitem{Mathlib2020}
The mathlib Community,
``The Lean mathematical library,''
\emph{Proceedings of the 9th ACM SIGPLAN International Conference on
Certified Programs and Proofs (CPP 2020)},
pp.~367--381, 2020.
\textsc{doi}: \href{https://doi.org/10.1145/3372885.3373824}%
{10.1145/3372885.3373824}.

\bibitem{MarklofStrombergsson2017}
Jens Marklof and Andreas Str\"{o}mbergsson,
``The three gap theorem and the space of lattices,''
\emph{The American Mathematical Monthly}, vol.~124, no.~8, pp.~741--745,
2017.

\bibitem{Moura2021}
Leonardo de Moura, Sebastian Ullrich,
``The Lean~4 theorem prover and programming language,''
\emph{Proceedings of the 28th International Conference on Automated
Deduction (CADE-28)},
LNAI, vol.~12699, pp.~625--635, Springer, 2021.
\textsc{doi}: \href{https://doi.org/10.1007/978-3-030-79876-5_37}%
{10.1007/978-3-030-79876-5\_37}.

\bibitem{Meyer1972}
Yves Meyer,
\emph{Algebraic Numbers and Harmonic Analysis},
North-Holland, Amsterdam, 1972.

\bibitem{Senechal2009}
Marjorie Senechal,
\emph{Quasicrystals and Geometry},
Cambridge University Press, 2009.
ISBN: 978-0-521-57541-6.

\bibitem{SpeechTao}
Terence Tao,
``Terence Tao on Yves Meyer's work on wavelets,''
video available at: \url{https://youtu.be/AnkinNVPjyw}, accessed
December 15, 2024.

\bibitem{ThreeGap}
``Three-gap theorem,''
Wikipedia,
\url{https://en.wikipedia.org/wiki/Three-gap_theorem}, accessed April 2025.

\end{thebibliography}

\end{document}
